% ---
% id: protocol-compiler/rem:block-restriction
% kind: rem
% title: "Restriction of the feasible set to block-diagonal states"
% summary: >-
%   The feasible set ranges over sector-block-diagonal states alone, a
%   restriction in the unsafe direction that the operator type carries. The
%   consuming engine's master reduction proves it sound by sector dephasing.
% prominence: standard
% status: accepted
% uses: [def:feasible]
% mentions: [meat-core/thm:reduction, meat-core/cor:reduction-constrained, meat-core/prop:a-neutral]
% assumes: []
% refs: []
% provenance: newly-derived
% lean:
% notes: "The condition is discharged by a theorem of the consuming engine and no claim declares it, so it is a remark and not an assumption."
% ---

\begin{remark}\label{rem:block-restriction}
\emph{The restriction.} The consuming engine's argument is sound when its
minimization ranges over the \emph{sector-block-diagonal} $\omega$ of
\eqref{eq:feas} rather than over all of $\Lin{A\Qsys}$: every state that
argument ranges over is carried into sector-block-diagonal form at no increase
of the bound. This module states the restriction and does not discharge it.
Its direction is the \emph{unsafe} one --- a smaller domain gives a larger
infimum, so an unlicensed restriction overstates a key length. And it is
carried by the operator \emph{type}: an $\omega$ \emph{is} a list of blocks
and no expressible object has an off-sector entry
(Definition~\ref{def:feasible}), so nothing verifies it and nothing can. The
argument that the sectors are a symmetry of every state the security proof
must admit belongs to the consuming engine and to the reductions module; this
module could not make it, since it never sees the reachable set.
Definition~\ref{def:feasible} \emph{defines} the \tref{feasible set}, and it
is not unsound if the restriction is unlicensed.

\emph{Where it is proved.} The consuming engine's master reduction
(\xref{meat-core/thm:reduction}, step~1 of its proof, the sector dephasing,
and \xref{meat-core/cor:reduction-constrained} for a non-empty list) proves
the restriction: an arbitrary feasible state is dephased, feasibility and
every constraint value survive because every $F_v$ and every $\hat O_k$ is
sector-block-diagonal, and the rest of the proof bounds the entropy by the
objective at the dephased state.

\emph{In the general format} (Definition~\ref{def:feasible}, the set of
\eqref{eq:feas-general}) the restriction is along the \emph{joint} sectors:
along the sectors of the \tref{measured system} as before and along Alice's
declared sectors. It is proved by the same step of the same theorem, which
dephases by the joint pinching, under preconditions that hold by
\emph{type}: Alice's measurement slot, $\tau_A$ and every right-hand side are
\tref[block list]{block lists} against a \emph{declared} factor tagged $A$,
as $\{F_v\}$ and every $\hat O_k$ are against the
\tref[measured factor]{measured factors} (\xref{meat-core/thm:reduction},
Step~1, with \xref{meat-core/prop:a-neutral} for the statement that the
restriction along Alice's sectors changes no value). The direction is
unchanged and so is the warning: the licence is that theorem under a declared
factor. A producer that emits a block-diagonal $\tau_A$ or $\{M_x\}$ without
a factor tagged $A$ has declared one Alice sector, and nothing downstream may
read a restriction off the shape of its data.
\end{remark}
